\section*{Resumen}

El objetivo de esta investigación es analizar la evolución de las disparidades
socioeconómicas y la convergencia regional en el Perú durante el periodo
2007--2019, considerando las posibles heterogeneidades espaciales existentes
entre los departamentos. El estudio comprende 24 departamentos y emplea dos
indicadores: el Índice de Desarrollo Humano y el producto bruto interno real
per cápita relativo. La estrategia empírica combina estadística descriptiva,
densidades kernel, análisis cartográfico, el índice global de Moran, regresiones
de convergencia beta mediante mínimos cuadrados ordinarios, modelos de error
espacial, especificaciones Durbin de error y regresiones geográficamente
ponderadas.

Los resultados evidencian comportamientos diferenciados. En el caso del
desarrollo humano, el coeficiente asociado con el nivel inicial es positivo y
estadísticamente significativo, por lo que no se identifica convergencia beta
durante el periodo analizado. Aunque la dispersión relativa del IDH disminuye,
los departamentos con mayores niveles iniciales registraron incrementos
superiores. En contraste, el coeficiente del PBI real per cápita relativo
inicial es negativo y altamente significativo, lo que revela un proceso de
convergencia económica regional. Los índices globales de Moran no muestran
autocorrelación espacial estadísticamente significativa; sin embargo, la
regresión geográficamente ponderada identifica diferencias territoriales en la
intensidad de la convergencia económica. En conjunto, los resultados sugieren
que la reducción de las brechas de ingreso no necesariamente estuvo acompañada
por una convergencia equivalente en las dimensiones del desarrollo humano.

\vspace{0.3cm}

\noindent\textbf{Palabras clave:} convergencia regional; desarrollo humano;
PBI per cápita; econometría espacial; heterogeneidad espacial; Perú.

\vspace{0.4cm}

\noindent\textbf{Clasificación JEL:} C21, O18, O47, R11, R12.
